% Bar-Ilan PhD Students Conference 2026 — MHM Pipeline
% Single source of truth: slide deck + Hebrew speaker notes.
%
% Build the slide deck (default):
%   cd docs/presentations
%   latexmk -xelatex bar-ilan-phd-pipeline-deck.tex
%
% Build the speaker-notes-only PDF:
%   Uncomment the \setbeameroption{show only notes} line below,
%   then run latexmk again. (Or maintain two thin wrappers that
%   \input this file.)
%
% Audit comments below each slide (% AUDIT:) record the dedupe /
% staleness fixes made when consolidating from the original HTML
% deck + bar-ilan-phd-pipeline-speaker-notes-he.tex.
%
% 2026-06-10: restructured from 12 slides into a three-act arc
% (10 slides). Distant supervision promoted to the Act-I hook;
% old "authority engine" (5) + "layers ranked" (7) merged into one
% slide; old data slide (3) folded into the numbers slide (now 6).
% FACT FIXES: design corpus "~50,000" had no source -> uses the
% verified 123,621 NLI count; ontology coverage "25% -> 87%" had no
% source for the 25% baseline -> now states "-> 87%" only.

\documentclass[aspectratio=169,12pt]{beamer}

\usepackage{fontspec}
\usepackage{polyglossia}
\setdefaultlanguage{english}
\setotherlanguage{hebrew}
\setmainfont{Times New Roman}
\newfontfamily\hebrewfont[Script=Hebrew,Scale=1.06]{Arial Unicode MS}

\usepackage{xcolor}
\usepackage{booktabs}
\usepackage{array}

% Theme — default Beamer. Swap to metropolis if installed.
\usetheme{default}
\setbeamertemplate{navigation symbols}{}
\setbeamerfont{frametitle}{size=\large,series=\bfseries}
\setbeamercolor{frametitle}{fg=black!85}
\setbeamercolor{normal text}{fg=black!85}
\setbeamercolor{itemize item}{fg=black!60}
\setbeamercolor{itemize subitem}{fg=black!60}

% Notes layout. Default = slides only. Uncomment one of the
% lines below to compile the notes-only / side-by-side variants.
% \setbeameroption{show only notes}
% \setbeameroption{show notes on second screen=right}
\setbeamertemplate{note page}{%
  \vspace{1ex}%
  \insertnote%
}

\title{Mapping Hebrew Manuscripts}
\subtitle{From MARC to Linked Open Data --- with the researcher in the loop}
\author{Alexander Goldberg}
\institute{Department of Information Science and Applied Artificial Intelligence \\ Bar-Ilan University}
\date{PhD Students Conference 2026}

\begin{document}

% ================================================================
% ACT I — THE PROBLEM & THE IDEA  (slides 1–3)
% ================================================================
% Narrative spine: a catalog full of knowledge, locked in a format
% built for shelving — and one idea (the catalog can supervise its
% own transformation) that unlocks it at scale.

% ─── Slide 1 — Title ────────────────────────────────────────────
\begin{frame}[t]{Mapping Hebrew Manuscripts}
  \vspace{0.4em}
  {\large From MARC to Linked Open Data --- with the researcher in the loop\par}
  \vspace{1.5em}
  \begin{columns}[T,onlytextwidth]
    \begin{column}{0.55\textwidth}
      \begin{tabular}{@{}lr@{}}
        \toprule
        MARC fields enriched against authority & \textbf{13} \\
        Authority sources & \textbf{4} \\
        Linked-data channels & \textbf{3} \\
        \bottomrule
      \end{tabular}
    \end{column}
    \begin{column}{0.42\textwidth}
      \textbf{Alexander Goldberg}\\
      Supervisors:\\
      Prof.\ Gila Prebor\\
      Dr.\ Avshalom Elmalech
    \end{column}
  \end{columns}
  \vfill
  {\footnotesize PhD Students Conference 2026 \quad\textbar\quad Bar-Ilan University}
\end{frame}
\note{\begin{hebrew}
שלום, אני אלכס גולדברג, מהמחלקה למדעי המידע ויישומי בינה מלאכותית, בהנחיית פרופ' גילה פריבור וד"ר אבשלום אלמליח.
\par
אני רוצה להתחיל בתמונה אחת. בספרייה הלאומית של ישראל יש מאות אלפי רשומות קטלוג של כתבי יד עבריים, ובתוכן מאות שנים של ידע --- מי כתב, מי העתיק, מי החזיק, ומאיפה הגיע כל כתב יד. אבל הידע הזה כלוא בפורמט שנבנה כדי לסדר ספרים על מדף, לא כדי לשאול עליו שאלות מחקר.
\par
ההרצאה הזאת היא סיפור של שחרור הידע הזה --- איך הופכים קטלוג ספרני לנתונים מקושרים פתוחים שאפשר באמת לחקור, ואיך עושים את זה כך שהחוקר נשאר במרכז. שלושת המספרים שעל השקף --- שלוש-עשרה שדות MARC, ארבעה מאגרי סמכות, שלושה ערוצי נתונים מקושרים --- הם תמצית המסע, ונחזור אליהם בסוף.
\end{hebrew}}

% ─── Slide 2 — The locked catalog (the problem) ────────────────
\begin{frame}[t]{A catalog full of knowledge --- locked}
  \begin{columns}[T,onlytextwidth]
    \begin{column}{0.48\textwidth}
      \textbf{What's in there}
      \begin{itemize}
        \item 123{,}621 NLI Hebrew-manuscript MARC records
        \item Authors, scribes, owners, collections
        \item Dates, places, works, genres
      \end{itemize}
    \end{column}
    \begin{column}{0.48\textwidth}
      \textbf{Why it stays locked}
      \begin{itemize}
        \item Names are local strings --- no stable authority ID
        \item Structured fields carry no external links
        \item Much of the story sits in free-text notes
      \end{itemize}
    \end{column}
  \end{columns}
  \vfill
  {\footnotesize MARC is stable cataloguer markup, built for shelving --- not for questions. $\rightarrow$ How do we unlock it, at scale, without armies of annotators?}
\end{frame}
\note{\begin{hebrew}
נתחיל מהבעיה. הפער המרכזי אינו שהמידע חסר --- הוא עשיר מאוד. הבעיה היא שהוא כלוא.
\par
בקטלוג של הספרייה הלאומית יש מאה עשרים ושלושה אלף שש מאות עשרים ואחת רשומות MARC של כתבי יד עבריים, ובהן מידע על כותבים, מעתיקים, בעלים, אוספים, תאריכים, מקומות, חיבורים וז'אנרים. אבל גם בשדות המבְנים עצמם, ערכים כמו שמות אנשים ומקומות הם בעיקר מחרוזות מקומיות --- בלי קישור אל מזהה יציב במאגרי סמכות בינלאומיים. חלק גדול מן הסיפור גם יושב בהערות חופשיות, בעברית, בארמית, ובניסוח שהשתנה לאורך עשרות שנות קטלוג.
\par
\textenglish{MARC} הוא פורמט מצוין לעבודה ספרנית, אבל הוא נבנה כדי לסדר ספרים על מדף --- לא כדי לשאול עליו שאלות. וכאן עולה השאלה שמניעה את כל המחקר: איך משחררים את הידע הזה, בקנה מידה גדול, \emph{בלי} צבא של אנשים שמתייגים רשומות ידנית?
\end{hebrew}}

% ─── Slide 3 — The insight: the catalog supervises itself ──────
% AUDIT: promoted distant supervision from a mid-talk aside to the
% Act-I hook. It is the single most interesting idea; it now opens
% the methodology rather than interrupting it.
\begin{frame}[t]{The idea: let the catalog teach the machine}
  \textbf{The catalog already annotated itself.}
  \vspace{0.6em}
  \begin{enumerate}
    \item A name sits in a \emph{structured} field --- \texttt{100}, \texttt{700}, \texttt{505}.
    \item The \emph{same} name reappears inside a \emph{free-text} note --- \texttt{500}, \texttt{561}.
    \item Structured value = the label; the note = the context. Together: one labelled training example.
  \end{enumerate}
  \vfill
  {\footnotesize \emph{Distant supervision} --- no manual annotation. Reusable: any MARC catalog with structured + free-text fields can train its own models.}
\end{frame}
\note{\begin{hebrew}
התשובה מתחילה ברעיון אחד, והוא לדעתי החלק הכי יפה במחקר: הקטלוג כבר תייג את עצמו.
\par
שימו לב למבנה הכפול של רשומת \textenglish{MARC}. שם של אדם יושב בשדה מובְנה --- למשל 100, 700, או 505. ואותו שם בדיוק מופיע שוב, הפעם בתוך הערה חופשית --- 500 או 561. הערך המובְנה נותן לנו את התווית; ההערה החופשית נותנת את ההקשר. יחד, השניים יוצרים דוגמת אימון מתויגת --- בלי שאף אדם תייג דבר. השיטה הזאת נקראת \textenglish{distant supervision}.
\par
זאת לא רק טכניקה נוחה, זאת תרומה מתודולוגית עצמאית: כל קטלוג \textenglish{MARC} בעולם שיש בו שדות מובְנים לצד שדות חופשיים יכול לאמן על עצמו מודלים, גם בלי תקציב לתיוג ידני. הרעיון הזה --- לתת לקטלוג ללמד את המכונה --- הוא שמניע את כל הפייפליין שנראה עכשיו.
\end{hebrew}}

% ================================================================
% ACT II — THE SYSTEM  (slides 4–6)
% ================================================================
% Narrative spine: the four-phase pipeline; where the value actually
% comes from (authority, honestly ranked); and the evidence.

% ─── Slide 4 — The pipeline: read · enrich · review · publish ──
\begin{frame}[t]{Read \textbullet{} Enrich \textbullet{} Review \textbullet{} Publish}
  \begin{itemize}
    \item \textbf{Read.} MARC coded fields \emph{and} free-text notes --- neither dropped.
    \item \textbf{Enrich.} People, places, works linked against Mazal, VIAF, Wikidata, KIMA. This is where most new claims are made; the trained models feed reliable entities into this phase.
    \item \textbf{Review.} The researcher meets every claim in a workbench --- sees it, compares it to the evidence, edits it, approves it.
    \item \textbf{Publish.} Three linked representations: full HMO RDF graph, a project Wikibase, and public Wikidata --- one consistent source.
  \end{itemize}
  \vfill
  {\footnotesize Between every phase: an editable JSON the researcher can correct before the next phase runs. Next: let's open up \emph{Enrich}.}
\end{frame}
\note{\begin{hebrew}
הרעיון הזה הופך לפייפליין של ארבעה מהלכים, וכל אחד מהם מקרב אותנו מהקטלוג הנעול אל נתונים מקושרים שאפשר לשאול עליהם.
\par
\emph{קריאה.} המערכת קוראת את שדות ה-\textenglish{MARC} המובְנים ואת ההערות החופשיות גם יחד --- בלי לאבד את האחד או את השני.
\par
\emph{העשרה.} כאן יושב הלב של התרומה. הפייפליין מקשר אנשים, מקומות וחיבורים מול ארבעה מאגרי סמכות --- \textenglish{Mazal} של הספרייה הלאומית, \textenglish{VIAF}, ויקינתונים ו-\textenglish{KIMA}. זהו המהלך שמייצר את רוב הטענות החדשות, והמודלים שאימנו ב-\textenglish{distant supervision} מזינים אליו ישויות אמינות.
\par
\emph{סקירה.} הקטלוגנים והחוקרים פוגשים את כל הטענות בסביבת עבודה גרפית: רואים, בודקים מול הראיה, עורכים, ומאשרים. זאת הסיבה שהפרויקט נבנה לחוקר, ולא רק כסקריפט.
\par
\emph{פרסום.} הנתונים יוצאים בשלושה ייצוגים מקושרים --- גרף ה-\textenglish{HMO} המלא, ויקיבייס פרויקטלי, וויקינתונים הציבורי.
\par
בין כל שני מהלכים יושב קובץ \textenglish{JSON} שהחוקר יכול לתקן לפני שהמהלך הבא רץ. עכשיו בואו נפתח את מהלך ההעשרה --- כי שם, באופן מפתיע, נולד רוב הידע החדש.
\end{hebrew}}

% ─── Slide 5 — Authority engine + honest ranking (MERGE) ───────
% AUDIT: merged the old "Authority carries the work" (Slide 5) and
% "Layers in order of contribution" (Slide 7) into ONE slide. The
% "NER is secondary" point is now made exactly ONCE — on the right.
\begin{frame}[t]{Where the value comes from --- honestly}
  \begin{columns}[T,onlytextwidth]
    \begin{column}{0.50\textwidth}
      \textbf{The engine: authority enrichment}
      \begin{itemize}
        \item \textbf{13 MARC fields $\rightarrow$ 4 authorities}\\
          {\footnotesize\texttt{100 110 111 600 610 651 700 710 711 752 800 810 811}}\\
          {\footnotesize Mazal (NLI) \textbullet{} VIAF \textbullet{} Wikidata \textbullet{} KIMA}
        \item \textbf{VIAF cluster harvest} --- one match $\rightarrow$ 2.3$\times$ extra IDs (GND, LCCN, ISNI, BnF)
        \item \textbf{Wikidata date back-fill} --- SPARQL on P569/P570
      \end{itemize}
    \end{column}
    \begin{column}{0.46\textwidth}
      \textbf{Ranked by measured contribution}
      \begin{tabular}{@{}p{0.05\textwidth}p{0.80\textwidth}@{}}
        \toprule
        1 & \textbf{Authorities} --- new-claim engine \\
        2 & \textbf{Rules} --- Hebrew dates, structure \\
        3 & \textbf{Classifiers} --- genre, routing \\
        4 & \textbf{NER} --- strong, net-new small \\
        \bottomrule
      \end{tabular}
    \end{column}
  \end{columns}
  \vfill
  {\footnotesize Order set by measurement, not by what looks impressive. The trained models are a safety net for the long tail --- not the primary engine.}
\end{frame}
\note{\begin{hebrew}
זה השקף שמסביר מאיפה באמת מגיעה התוספת --- ואני רוצה להיות כן לגביו, כי התשובה אינטואיטיבית פחות ממה שמצפים.
\par
מצד שמאל, \emph{המנוע}: העשרת סמכות. הפייפליין מזין שלוש-עשרה תגיות \textenglish{MARC} אל מאגרי הסמכות --- לא רק 100/700 הקלאסיים, אלא גם 600/610 לנושאים, 651 ו-752 למקומות, ו-800/810/811 לכניסות סדרה --- ומזהה כל ישות מול ארבעה מאגרים משלימים: \textenglish{Mazal}, \textenglish{VIAF}, ויקינתונים, ו-\textenglish{KIMA}. שני שיפורים מכפילים את הערך: \emph{קציר אשכול \textenglish{VIAF}}, שבו התאמה אחת מביאה איתה פי 2.3 מזהים בינלאומיים נוספים --- \textenglish{GND, LCCN, ISNI, BnF}; ו\emph{השלמת תאריכים מוויקינתונים} דרך \textenglish{P569/P570}, כשהמקורות האחרים שותקים.
\par
מצד ימין, הדירוג הכן של ארבע השכבות --- לפי מה שהמדידות הראו, ולא לפי מה שמרשים על הנייר. ראשונה הסמכות; שנייה החוקים, שחזקים במיוחד במקום שבו המבנה יציב, כמו תאריכים; שלישית המסווגים, שמשלימים שדות חסרים; ורק רביעית --- מודלי ה-\textenglish{NER}.
\par
וכאן הנקודה הכנה, ואומר אותה פעם אחת בבירור: מודלי ה-\textenglish{NER} שלנו \emph{טובים} במשימת החילוץ, אבל אחרי שלב הסמכות והחוקים, תוספת הטענות הנקייה שלהם בקורפוס הזה קטנה. הם רשת ביטחון לזנב הארוך --- לרשומות שבהן הישות חיה רק בהערה חופשית --- ולא המנוע המרכזי. בשקף הבא נראה את המספרים המדויקים שמאחורי כל זה.
\end{hebrew}}

% ─── Slide 6 — The numbers (evidence) ──────────────────────────
% AUDIT: the 68-record evaluation set (was its own Slide 3) is now
% folded in here as the measurement context for the numbers.
% AUDIT: FACT FIX — "25% → 87%" had no source for the 25% baseline;
% changed to "→ 87%" (the 87% is verified; the baseline is not).
\begin{frame}[t]{What each layer contributes --- in numbers}
  \begin{tabular}{@{}p{0.28\textwidth}p{0.22\textwidth}p{0.44\textwidth}@{}}
    \toprule
    \textbf{Layer} & \textbf{Headline} & \textbf{What it means} \\
    \midrule
    Baseline & $\sim$8 claims/MS & direct MARC conversion, no extras \\
    Authority & +2.3$\times$ IDs & VIAF cluster harvest per matched person \\
    Rules & 22\% $\rightarrow$ 97\% & manuscripts with a parsed creation date \\
    Classifiers & 28 $\rightarrow$ 39 & records with a genre (of 68) \\
    NER + cls. & 80.4\% F1 strict & person v3 (name+role); net new small \\
    LOD bridge & $\rightarrow$ 87\% & HMO classes reachable via Wikidata; 28.1 claims/MS \\
    \bottomrule
  \end{tabular}
  \vfill
  {\footnotesize Measured on a 68-record evaluation set (mostly a 17th-century sample). Baseline $\rightarrow$ full pipeline: 8.0 $\rightarrow$ 28.1 claims/MS ($3.5\times$). Person v3: 80.4\% strict / 86.8\% name-only / 92.6\% role acc. Provenance NER 95.91\%; Contents NER 99.99\% on the 2 records where MARC 505 appears.}
\end{frame}
\note{\begin{hebrew}
לפני המספרים, מילה על מה שמדדנו: מערך הערכה קבוע של 68 רשומות, רובן מדגם של כתבי יד מן המאה ה-17. מערך קטן מספיק כדי לבדוק כל רשומה ביד, ומגוון מספיק כדי להיות כן. כל מספר שאני אומר עכשיו חוזר למערך הזה.
\par
\emph{קו הבסיס.} המרה ישירה של \textenglish{MARC}, בלי שום שכבה נוספת, נותנת בערך שמונה טענות לכתב יד. הפייפליין המלא מגיע ל-28.1 --- שיפור של פי 3.5.
\par
\emph{סמכות --- המנוע.} כל אדם שמתאים ל-\textenglish{VIAF} מקבל פי 2.3 מזהים בינלאומיים מן האשכול. זאת השכבה שהופכת כל ישות ממחרוזת לישות מקושרת.
\par
\emph{חוקים.} כיסוי תאריך היצירה עולה מ-22 אחוז ל-97 אחוז --- כמעט כולו מחוקים שמפענחים תאריכים מקודדים וניסוחי מאות בעברית.
\par
\emph{מסווגים.} כיסוי הז'אנרים עולה מ-28 מתוך 68 רשומות ל-39 מתוך 68.
\par
\emph{NER --- מדוד ושקוף.} מול תוויות הזהב, מודל האנשים מגיע ל-80.4 אחוז \textenglish{F1} בהתאמה קפדנית של שם ותפקיד, 86.8 אחוז בשם בלבד, ו-92.6 אחוז דיוק תפקיד. הפרובננס מגיע ל-95.91 אחוז, והתוכן ל-99.99 אחוז --- אבל שימו לב, רק בשתי רשומות מתוך 68 בכלל יש שדה 505. המודלים טובים; פשוט עיקר התרומה כבר נתפס על ידי הסמכות והחוקים.
\par
\emph{הגשר.} כל השכבות מזינות את ההקרנה לוויקינתונים, שמגיעה לכ-87 אחוז ממחלקות האונטולוגיה ו-28.1 טענות בממוצע לכתב יד --- ואת זה נראה בשקף הבא.
\end{hebrew}}

% ================================================================
% ACT III — THE HUMAN & THE PAYOFF  (slides 7–10)
% ================================================================
% Narrative spine: trust comes from the human in the loop; the output
% is three layers of LOD; the payoff is new research questions; close
% on the four contributions.

% ─── Slide 7 — Wikidata Studio + the eval-agent ────────────────
\begin{frame}[t]{Why you can trust the numbers: the human in the loop}
  \textbf{Wikidata Studio --- the workbench}
  \begin{itemize}
    \item Three-axis review: \emph{source} (which MARC field or model), \emph{entity type}, \emph{role}.
    \item Biodata-comparison dialog --- dates, places, name variants, occupations side by side from MARC, Mazal, VIAF, KIMA.
    \item Confidence tooltip --- which sources matched, which guards fired, which years were pulled.
    \item 14 pre-save checks, resumable review state, paging at 25/50/100 rows.
  \end{itemize}
  \textbf{The eval-agent}
  \begin{itemize}
    \item A Gemini-based judge audits \emph{both} entity extraction \emph{and} authority matches.
    \item One shared approval across every editor --- approve once, reflected everywhere.
    \item \textbf{Auto-fix}: a high-confidence correction appears inline --- text is updated, but approval stays human.
  \end{itemize}
  \vfill
  {\footnotesize The agent suggests; the human decides. Every number on the previous slide passed through this.}
\end{frame}
\note{\begin{hebrew}
ועכשיו, הצד שהופך את כל המספרים האלה לאמינים. כי אפשר לחלץ ולהעשיר כמה שרוצים --- אם אף אדם לא בדק את התוצאה, חוקר לא יוכל לסמוך עליה. \textenglish{Wikidata Studio} היא סביבת העבודה שמחזירה את האדם למרכז.
\par
\emph{סקירה בשלושה צירים.} כל שורה היא ישות שראויה לסקירה: כל ערך \textenglish{MARC} שעבר העשרה, כל ישות שזיהה מודל \textenglish{NER}, כל מקום שזיהה \textenglish{KIMA}. אפשר לסנן לפי \textbf{מקור} --- איזה שדה או איזה מודל --- לפי \textbf{סוג ישות}, ולפי \textbf{תפקיד}. שלושה צירים שהם בעצם השפה של הקטלוגנית.
\par
\emph{דיאלוג השוואת ביוגרפיה.} ההחלטה אם להאמין להתאמת סמכות אינה נופלת על שם בלבד, אלא על תאריכים, מקומות, וריאנטים בשמות, ועיסוקים. לחיצה פותחת מסך שמראה את כולם זה לצד זה --- מ-\textenglish{MARC}, \textenglish{Mazal}, \textenglish{VIAF}, ו-\textenglish{KIMA} --- עם הבלטה צבעונית של מה תואם ומה שונה. מנרמלים את הטקסט ומסירים תוויות כמו ``בן'' או ``\textenglish{de}'', כדי לא להישבר על דמיון שטחי.
\par
\emph{\textenglish{Tooltip} של ביטחון.} ציון \textenglish{HIGH} או \textenglish{LOW} אינו תווית סתומה: ריחוף עכבר מראה מי תאם, אילו שומרי בקרה ירו, ואילו שנים ביוגרפיות נשלפו. ו-14 בדיקות לפני שמירה חוסמות פריטים פגומים. מצב הסקירה נשמר, אפשר לעצור ולחזור, ודפדוף של 25/50/100 שורות שומר על מהירות גם בקורפוסים גדולים.
\par
\emph{סוכן ההערכה.} מעל הסקירה הידנית, סוכן מבוסס \textenglish{Gemini} מבקר את הפלט בשתי נקודות --- \emph{גם את חילוץ הישויות וגם את התאמות הסמכות} --- עם אישור משותף בכל המסכים: אישור פעם אחת תקף בכולם. וכשהסוכן מזהה ישות נכונה חלקית, למשל שם שחסר בו פטרונים, הוא מציע \textbf{תיקון אוטומטי} ברמת ביטחון גבוהה --- הלחצן מעדכן את הטקסט בלבד, וההחלטה לאשר נשארת בידי האדם. הסוכן מציע; האדם מכריע. כל מספר שראינו בשקף הקודם עבר דרך התהליך הזה.
\end{hebrew}}

% ─── Slide 8 — MARC → LOD in three layers ──────────────────────
\begin{frame}[t]{The output: Linked Open Data, in three layers}
  \begin{itemize}
    \item \textbf{HMO graph (RDF).} The full scholarly graph --- manuscripts, works, persons, places, copying and ownership events, supporting evidence. Now persisted in PostgreSQL, always available.
    \item \textbf{Project Wikibase} (\texttt{mhm-hmo.wikibase.cloud}). The same graph in a familiar browse-and-edit interface --- no ontology class dropped.
    \item \textbf{Public Wikidata.} A projection: after all model improvements, $\sim$87\% of the ontology's classes are reachable via Wikidata, $\sim$28.1 claims per manuscript.
  \end{itemize}
  \vfill
  {\footnotesize One source, three audiences --- domain experts, cataloguers, the open web. All three are queryable from the \textbf{Linked Data Explorer} tab.}
\end{frame}
\note{\begin{hebrew}
עכשיו שכל טענה נבדקה, אפשר לפרסם --- והפרסום עצמו הוא תרומה. מאותו צינור עיבוד יוצאים \emph{שלושה} ייצוגים של נתונים מקושרים, כל אחד לקהל אחר.
\par
\emph{הראשון, גרף ה-\textenglish{HMO}.} הייצוג המחקרי המלא, מבוסס על האונטולוגיה שלנו לכתבי יד עבריים: כתב היד, החיבורים, האנשים, המקומות, אירועי העתקה ובעלות, והראיות שעליהן הטענות נשענות. החידוש: הגרף נשמר עכשיו ב-\textenglish{PostgreSQL}, כך שהוא זמין לאחר כל הפעלה מחדש של השרת, בלי \textenglish{rebuild}.
\par
\emph{השני, ויקיבייס פרויקטלי.} אותו גרף נטען לויקיבייס ייעודי בכתובת \texttt{mhm-hmo.wikibase.cloud}, כך שאפשר לדפדף, לחפש ולערוך את כל הישויות בממשק מוכר --- בלי לאבד אף קלאסה מן המבנה המחקרי.
\par
\emph{השלישי, ויקינתונים הציבורי.} כאן מתבצעת הקרנה. אחרי כל שיפורי המודל, כ-87 אחוז ממחלקות האונטולוגיה ניתנות להגעה דרך ויקינתונים, ובממוצע 28.1 טענות לכתב יד. הופכים את ה-\textenglish{MARC} לנתונים מקושרים פתוחים --- בלי לאבד לא את העומק המחקרי ולא את הנגישות הציבורית. ושלושת הייצוגים האלה ניתנים לשאילתה ישירה, כפי שנראה עכשיו.
\end{hebrew}}

% ─── Slide 9 — New questions + Linked Data Explorer (payoff) ───
\begin{frame}[t]{The payoff: questions you couldn't ask before}
  \begin{itemize}
    \item Which works \emph{travel together} in the same manuscripts? --- textual clusters from the corpus itself.
    \item Whose name \emph{reappears across collections}? --- the social network of manuscript production and transmission.
    \item Whose \emph{ownership chain} crosses which libraries, families, and places, across centuries?
    \item How did manuscripts \emph{move} between places and periods?
  \end{itemize}
  \vfill
  {\footnotesize \textbf{Linked Data Explorer} tab: a three-layer SPARQL console (HMO graph / Wikibase / Wikidata) with 8 pre-built scholarly templates. A provenance header confirms every entity shown was human-reviewed in Wikidata Studio.}
\end{frame}
\note{\begin{hebrew}
וזה הרגע שבשבילו עשינו את כל הדרך. ברגע שהפלט הוא נתונים מקושרים פתוחים --- ולא רק רשומות \textenglish{MARC} --- אפשר לשאול שאלות שכמעט בלתי-אפשרי לשאול מן הקטלוג המקורי.
\par
אילו \emph{חיבורים מופיעים יחד} באותם כתבי יד, וכך לזהות אשכולות טקסטואליים מן הקורפוס עצמו. אילו \emph{אנשים} --- מעתיקים, בעלים, נמענים --- חוזרים על פני כמה אוספים, וכך לשרטט את הרשת החברתית של ייצור והעברה. אילו \emph{שרשראות בעלות} חוצות ספריות, משפחות ומקומות לאורך מאות שנים. ואיך כתבי יד \emph{נעו} בין מקומות ותקופות.
\par
\textbf{וכל זה לא תיאורטי.} לשונית ה-\textenglish{Linked Data Explorer} מציגה קונסולת \textenglish{SPARQL} בשלוש שכבות --- גרף ה-\textenglish{HMO} בתוך המערכת, ויקיבייס פרויקטלי, וויקינתונים הציבורי --- עם שמונה תבניות שאילתה מוכנות מראש לאשכולות חיבורים, רשת מעתיקים, שרשראות בעלות, ופיזור גיאוגרפי. בראש הלשונית יושב כותר מקורות עם תג אמינות, שמאשר שכל ישות שמוצגת עברה סקירת אדם ב-\textenglish{Wikidata Studio}. חוקר שמריץ שאילתה יודע שהקישורים מאחורי כל מזהה נבחנו --- ולא נוחשו. הפייפליין איננו סוף המחקר; הוא התשתית שמאפשרת מחקר, וכלי המחקר מובנה בתוכו.
\end{hebrew}}

% ─── Slide 10 — Close: four contributions ──────────────────────
\begin{frame}[t]{Four contributions, one pipeline}
  \begin{enumerate}
    \item \textbf{Authority enrichment as the primary engine.}\\
      13 MARC fields against 4 authorities, with VIAF cluster harvest and Wikidata date back-fill.
    \item \textbf{LOD publication in three layers + Linked Data Explorer.}\\
      HMO RDF (persisted in PostgreSQL), project Wikibase, Wikidata --- one pipeline, three audiences, one live SPARQL console.
    \item \textbf{Wikidata Studio --- a research instrument.}\\
      Three-axis review, biodata comparison, 14 pre-save checks, resumable state.
    \item \textbf{The eval-agent --- extraction, authority audit, and Auto-fix.}\\
      An agentic judge over whole corpora; high-confidence NER fixes surfaced inline; one shared approval across every editor.
  \end{enumerate}
  \vfill
  \begin{center}
    \emph{The pipeline is not the end of the research --- it is the infrastructure that makes research possible.}\\[0.4em]
    {\footnotesize Live curator interface: \texttt{mhm-pipeline.org} \quad\textbar\quad \texttt{shvedbook@gmail.com}}
  \end{center}
\end{frame}
\note{\begin{hebrew}
אסכם, ואחזור לשלושת המספרים מן ההתחלה. ארבע תרומות, פייפליין אחד.
\par
\emph{ראשית --- העשרת סמכות כמנוע העיקרי.} שלוש-עשרה שדות \textenglish{MARC} מול ארבעה מאגרים, עם קציר אשכול \textenglish{VIAF} והשלמת תאריכים מוויקינתונים. זה מה שמייצר את רוב הטענות החדשות.
\par
\emph{שנית --- פרסום בשלושה ייצוגים, ועכשיו גם ניתן לשאילתה ישירה.} גרף \textenglish{HMO} מלא לחוקר (שמור ב-\textenglish{PostgreSQL}), ויקיבייס לדפדוף ועריכה, וויקינתונים לציבור --- ולשונית ה-\textenglish{Linked Data Explorer} שהופכת את השאלות המחקריות מתיאוריה למשהו שמריצים.
\par
\emph{שלישית --- \textenglish{Wikidata Studio} ככלי מחקר.} סקירה בשלושה צירים, השוואת ביוגרפיה, 14 בדיקות, ומצב סקירה שאפשר לעצור ולחזור אליו --- מה שמאפשר לחוקר לסמוך על כל טענה.
\par
\emph{רביעית --- סוכן ההערכה.} שכבת ביקורת אגנטית על חילוץ \emph{ועל} סמכות, עם תיקון אוטומטי בביטחון גבוה ואישור משותף. מודלי ה-\textenglish{NER} נשארים כתרומה מתודולוגית --- \textenglish{distant supervision} שניתן ליישום בכל קטלוג דומה --- ותרומתם נמדדה כאן במפורש.
\par
מבחינתנו, השאלה אינה רק איך מעבירים נתונים מפורמט אחד לאחר. השאלה היא איך הופכים ידע קטלוגי, שנבנה במשך שנים על ידי מומחים, לתשתית מחקרית שאפשר לשאול עליה שאלות חדשות --- בצורה שקופה, זהירה, ואחראית, כשהאדם נשאר במרכז. הכלי פועל היום כסביבת ענן ב-\textenglish{mhm-pipeline.org}, ומאפשר לצוות לסקור, לערוך ולאשר ישויות ביחד. תודה רבה.
\end{hebrew}}

\end{document}
